% !TEX root = ../Tesis.tex
\chapter{Marco teórico}
\section{Antecedentes}
Un estudio realizado por Gorski et al., 2015 en Salt Lake City
determinó que el vapor de agua proveniente de la combustión de
combustibles fósiles (VDC) podía ser identificada a través de la
medición de isótopos estables.

El grupo de investigación tomó muestras de vapor ambiental bombeada a
través de una tubería de cobre hasta un vaporizador acoplado a un
espectrómetro Picarro L2130-i CRDS. El muestreo se realizó durante el
invierno de diciembre de 2013 a enero de 2014 (Gorski et al., 2015).

De acuerdo con el análisis isotópico, se determinó que el agua asociada
a la combustión presentaba valores de d-excess altamente negativos
(entre $-471$\,\textperthousand\ y $-183$\,\textperthousand). Al mismo
tiempo, durante episodios de inversión térmica, las concentraciones de
CO\textsubscript{2} aumentan y los valores de d-excess del vapor de
agua disminuyeron drásticamente (hasta $-17.2$\,\textperthousand). Esta
relación inversa ($r^{2} = 0.33$ a $0.64$) se observó en escalas de
varios días y en ciclos diurnos (Gorski et al., 2015).

A través de un modelo de mezcla de dos fuentes (vapor de combustión y
vapor de la atmósfera libre), se estimó que el agua de combustión
representó hasta el 13\% del vapor total en la capa límite durante este
periodo. Las concentraciones de este vapor alcanzaron hasta 200 ppm en
el aire (Gorski et al., 2015).

Asimismo, la investigación realizada por Y. Yang y Yoshimura (2024)
destaca la correlación negativa entre el valor de d-excess del vapor
atmosférico y la concentración de CO\textsubscript{2} previamente
descrita por Gorski y colaboradores. El propósito del estudio fue
evaluar si el modelo IsoRSM mejoraba la precisión de las simulaciones
isotópicas al incorporar un modelo de emisión de VDC (Y. Yang y
Yoshimura, 2024).

La investigación implementó dos análisis independientes con un mes de
duración con datos obtenidos de bases como ERA5, la Red de observación
in situ MesoWest y la Base de Datos de Isótopos de Vapor de Agua
Estable. El primero, en Salt Lake City, realizado en enero de 2017. Se
situó en una región caracterizada por topografía de cuenca en una
meseta, donde las inversiones térmicas invernales generan condiciones
de atmósfera estable. El segundo se desarrolló en Beijing en enero de
2007. En este caso se realizó en una llanura sometida a la influencia
del monzón invernal (Y. Yang y Yoshimura, 2024).

Por otro lado, la composición isotópica del VDC se fijó como
$-134.4$\,\textperthousand, $9$\,\textperthousand\ y
$-206.4$\,\textperthousand\ para $\delta^{2}$H, $\delta^{18}$O y
d-excess, respectivamente, basados en experimentos realizados en Xi'an,
China. Para la estimación de la tasa de emisión de VDC a añadir se
consideró la tasa de emisión de CO\textsubscript{2} en zonas urbanas.
Se evaluaron tasas de $0{,}000$ hasta $0{,}003$\,g/m\textsuperscript{2}·s
de VDC para comprobar cuál presentaba la mejor correlación (Y. Yang y
Yoshimura, 2024).

Con ello, obtuvieron que la incorporación de VDC redujo
significativamente el d-excess del vapor de agua, especialmente cuando
la capa límite planetaria era estable y la humedad era baja. En Salt
Lake City hubo una disminución aproximadamente de
$5$\,\textperthousand, mientras que en Beijing fue más pronunciada
(Y. Yang y Yoshimura, 2024).

Mediante el método de trazadores de humedad, se estimó que la fracción
mensual de CDV era del 3.4\% en SLC y 4.7\% en Beijing. La fracción
calculada a partir de la ecuación de mezcla de dos extremos
(f\textsubscript{CDV\_iso}) mostró valores ligeramente inferiores:
2.3\% y 3.3\%, respectivamente.

\section{La atmósfera urbana}
La atmósfera terrestre es una fina capa de gases que rodea el planeta.
Esta capa es fundamental para la vida en la Tierra, ya que regula la
temperatura de su superficie y retiene el oxígeno y el agua, esenciales
para la vida (Borduas \& Donahue, 2018).

% IMAGEN REDISEÑADA
%\begin{figure}[H]
    \centering
    \includegraphics[width=0.8\textwidth]{figuras/nombre_imagen1.png}
    \caption{Descripción de la imagen.}
    \label{fig:atmosfera}
%\end{figure}

La capa límite planetaria (CLP) es la capa baja y esencialmente
turbulenta de la atmósfera. Responde a los cambios espaciales o
temporales en las propiedades de la superficie terrestre o a los
impactos procedentes de ella. La CLP crea una interfaz entre la
superficie y la tropósfera libre que se extiende hasta la tropopausa.
Una de las características de la CLP es que en la parte superior su
turbulencia es nula (Baklanov, 2024; Sokhi et al., 2024).

% IMAGEN REDISEÑADA
%\begin{figure}[H]
    \centering
    \includegraphics[width=0.8\textwidth]{figuras/nombre_imagen2.png}
    \caption{Descripción de la imagen.}
    \label{fig:clp}
%\end{figure}

La capa límite urbana (CLU) se define como la parte de la capa límite
planetaria que es impactada por la presencia de urbanización en la
parte baja. La CLU se extiende desde el nivel de los techos de los
edificios hasta una altura en la que los efectos de la superficie
urbana ya no son apreciables. Sin embargo, si la extensión de la
urbanización es lo suficientemente grande, la CLU puede abarcar toda la
profundidad de la CLP (Berry, 2022; Sokhi et al., 2024; Vallero, 2025).

La urbanización y las actividades antropogénicas pueden influir en
características atmosféricas a micro y mesoescala, por ejemplo
(Baklanov, 2024):

\begin{itemize}
    \item La heterogeneidad de las construcciones y otros elementos
    rugosos de la superficie afectan el régimen de turbulencia del
    flujo.
    
    \item El uso excesivo de materiales impermeables y la reducción de
    fracciones permeables alteran el régimen hidrometeorológico.
    
    \item La deposición de contaminantes.
    
    \item La emisión de flujos de calor antropogénicos derivados de las
    actividades humanas afecta al régimen térmico.
    
    \item La emisión de contaminantes como aerosoles, que impactan en la
    transferencia de radiación y la formación de nubes y precipitación.
    
    \item Los cañones urbanos (debido a la geometría de las
    construcciones) impactan en el régimen del flujo e intercambio de
    calor entre superficies (caminos y paredes).
\end{itemize}

Además, la CLU es más compleja que la capa límite rural (CLR). Durante
el día, a mesoescala, la CLU es cálida, turbulenta y contaminada. Está
formada por una capa superficial (CS) y una capa de mezcla (FIG). Con
un viento promedio, la CLU forma una pluma de advección sobre la CLR
situada a sotavento. Por otro lado, cuando no hay viento, se forma una
capa límite superficial urbana con forma de cúpula. En escala local, la
capa superficial es comprimida por la capa canopia urbana (CCU, límite
superior marcado por $Z_H$) situada entre sus edificios, que constituye
la mitad inferior de la subcapa de rugosidad (SCR, límite superior en
$Z_R$). La subcapa más alta de CS es la subcapa inercial (SCI, límite
superior en $Z_L$). La CLU suele estar limitada por una zona de
arrastre (ZA), por encima de la cual se encuentra la atmósfera libre
(AL). La ZA suele estar centrada en la base de una inversión de
cubierta elevada (IC) (Baklanov, 2024; Vallero, 2025).

Por otro lado, durante la noche, la capa de inversión superficial
nocturna (CISN) de la zona rural situada a barlovento es arrastrada
hacia la ciudad, lo que da lugar a una inversión elevada débil y poco
profunda. Suele quedar cubierta sucesivamente por la capa de rugosidad
nocturna (CR), a veces una IC, y luego por la AL (Baklanov, 2024;
Vallero, 2025). La capa límite nocturna es estable. Por ello, al emitir
contaminantes atmosféricos durante la noche, estos se dispersan
horizontalmente en capas finas, ya que apenas hay mezcla vertical
(Sokhi et al., 2024).

% IMAGEN REDISEÑADA
%\begin{figure}[H]
    \centering
    \includegraphics[width=0.8\textwidth]{figuras/nombre_imagen3.png}
    \caption{Descripción de la imagen.}
    \label{fig:clu_nocturna}
%\end{figure}

Al aumentar la velocidad del viento, aumenta la mezcla vertical de las
masas de aire. La estabilidad de estas masas toma en cuenta la
estabilidad estática y dinámica. La estabilidad estática se ve afectada
por las fuerzas de flotabilidad derivadas de las diferencias de
densidades. Mientras tanto, la estabilidad dinámica se basa en las
fuerzas de cizalladura (Vallero, 2025).

% IMAGEN REDISEÑADA
%\begin{figure}[H]
    \centering
    \includegraphics[width=0.8\textwidth]{figuras/nombre_imagen4.png}
    \caption{Descripción de la imagen.}
    \label{fig:estabilidad}
%\end{figure}